\section{Trusted base}
\label{sec:mechanization-03}

\begin{definition}[Trusted base]
For a particular checked result, the \emph{trusted base} is the set of
components and assumptions whose correctness is required but is not itself
established by the checked proof term.
\end{definition}

The trusted base normally includes the kernel executable and hardware.  It may
also include axioms, foreign code, an external solver, generated proof
certificates, or an unverified parser, depending on the workflow.  The trusted
base is result-specific; naming the proof assistant is not a complete trust
report.

\begin{method}[Mechanized-result report]
Report separately:
\begin{enumerate}
  \item the natural-language theorem intended;
  \item the exact formal theorem statement;
  \item the correspondence argument between them;
  \item the assistant and exact version;
  \item the imported modules;
  \item the axioms printed or otherwise inspected;
  \item the kernel-check result;
  \item external tools and certificate checking; and
  \item remaining human review of definitions and fidelity.
\end{enumerate}
\end{method}

